\section{Detailed Derivations of the Expansion Errors}
\label{appendix_expansion_error_detail}

Here we derive the expansion errors, $H_M^{\{\sech\}}$, $I_M^{\{\sech\}}$,
$H_M^{\{\tanh\}}$, $I_M^{\{\tanh\}}$, $H_M^{\{\sech^2\}}$, $I_M^{\{\sech^2\}}$ in detail.
We find that the integrals involved are resistant to direct solution by computer
algebra systems and are not readily solved via existing integration rules.  Our approach focuses
on breaking the integrals for the exact errors $H_M^{\{\varphi\}}$ apart into factors where
some factors are responsible for divergence and other factors are responsible for the oscillation
within the region of convergence, helping to reveal the closed form approximation for the
Gibb's like errors within the region of convergence, which we denote $I_M^{\{\varphi\}}$.

As an intermediate step in solving for the error in each expansion,
we reformulate the problem in terms of the general Meijer-G function which facilitates further
reduction via Mathematica \cite{Mathematica}.
In reducing and manipulating the error integrals into their useful factorizations
and in extracting their closed form approximations, we make use of identities from
Prudnikov's \foreignlanguage{russian}{Интегралы и Ряды Том 3, Специальные функции 
	Дополнительные Главы} (Integrals and Series Volume 3, More Special Functions) \cite{prudnikov_integrals_2003}.

\subsection{Hyperbolic Secant}
We begin with the exact error formula as derived in Chapter \ref{chapter_improved_expansions}, $H_M^{\{\sech\}}$, which follows directly from application of
(\ref{eq_general_remainder}):
\begin{align}
	\begin{split}
		&H_M^{\{\sech\}}(v) = \int_\sigma^{\infty} \!
		\sech\left(\frac{\pi}{2}\xi\right)
		\cos(\xi v) \mathrm{d}\xi
	\end{split}
	\\
	\begin{split}
		& \quad\quad + \frac{(-1)^{M+1}v^{2M+2}}{(2M+2)!} \int_0^{\sigma} \!
		\xi^{2M+2} \sech\left(\frac{\pi}{2}\xi\right)
		{}_{1}F_{2}\left(1; M+2, M+\frac{3}{2}; -\frac{v^2 \xi^2}{4}\right) \mathrm{d}\xi.
	\end{split}
\end{align}

As noted in Chapter \ref{chapter_improved_expansions}, the second term dominates
for large $v$ due to the $v^{2M+2}$ factor, causing the expected divergence
for $|v| \geq V$.  Meanwhile, the first term dominates for small $v$ leading
to Gibbs-like oscillation of the error for $|v| < V$.  Noting this effect,
the integral of (\ref{eq_tanh_remainder2}) can be solved exactly to analyze the
error, $I_M^{\{\sech\}}(v)$, inside the region of convergence for
an $M$ term expansion.  We begin to attack this integral by breaking up the integrand via
\begin{align}
	\sech\left(\frac{\pi \xi}{2}\right)\cos(\xi v) = 
	\frac{e^{\frac{\pi\xi}{2}} e^{-iv\xi}}{1 + e^{\pi \xi}}
	+
	\frac{e^{\frac{\pi\xi}{2}} e^{iv\xi}}{1 + e^{\pi \xi}}.
\end{align}
Integrating each term yields
\begin{align}
	\int_\sigma^\infty \! \frac{e^{\frac{\pi\xi}{2}} e^{\pm iv\xi}}{1 + e^{\pi \xi}} \mathrm{d}\xi
	= \sech(v) - \frac{2 e^{\frac{1}{2}\sigma(\pi \pm 2iv)}}{\pi \pm 2iv}
	{}_{2}F_{1}\left(1, \frac{1}{2} \pm \frac{iv}{\pi}; \frac{3}{2} \pm \frac{iv}{\pi}; -e^{\pi\sigma}\right).
\end{align}
This results in
\begin{align}
	I_M^{\{\sech\}}(v) &= 2 \sech(v) - \frac{1}{\pi}\left(u + u^*\right), \\
	u &= \frac{(e^{\pi\sigma})^{\frac{1}{2} + \frac{iv}{\pi}}}{\frac{1}{2} + \frac{iv}{\pi}}
	{}_{2}F_{1}\left(1, \frac{1}{2} + \frac{iv}{\pi}; \frac{3}{2} + \frac{iv}{\pi}; -e^{\pi \sigma}\right).
\end{align}
This can be equivalently expressed in terms of the Lerch Transcendent, $\Phi(z, a, b)$, by noting
from \cite{prudnikov_integrals_2003} that
\begin{align}
	{}_{2}F_{1}\left(1, b; b+1; z\right) = b \Phi(z, 1, b)
\end{align}
implies
\begin{gather}
	\begin{split}
		&I_M^{\{\sech\}}(v) = \int_\sigma^{\infty} \!
		\sech\left(\frac{\pi}{2}\xi\right)
		\cos(\xi v) \mathrm{d}\xi \\
		\ \ &= 2 \sech(v)
		- \frac{1}{\pi}\left(
		\left( e^{\pi\sigma}\right)^{\frac{1}{2} + \frac{iv}{\pi}}
		\Phi\left(-e^{\pi \sigma}, 1, \frac{1}{2} + \frac{iv}{\pi}\right)
		+
		\left( e^{\pi\sigma}\right)^{\frac{1}{2} - \frac{iv}{\pi}}
		\Phi\left(-e^{\pi \sigma}, 1, \frac{1}{2} - \frac{iv}{\pi}\right)
		\right).
	\end{split}
\end{gather}
Noting that the identity
\begin{align}
	\Phi\left(1-z, 1, a\right)
	= \frac{1}{\Gamma(a)}
	G_{2,2}^{2,2}
	\left(
	\begin{matrix}
		0, & 1-a \\
		0, & 0
	\end{matrix}\;\middle|\;z
	\right)
\end{align}
from \cite{prudnikov_integrals_2003} implies
\begin{align}
	\Phi\left(-e^{\pi\sigma}, 1, \frac{1}{2} \pm \frac{iv}{\pi}\right)
	= \frac{1}{\Gamma\left(\frac{1}{2} \pm \frac{iv}{\pi}\right)}
	G_{2,2}^{2,2}
	\left(
	\begin{matrix}
		0, & \frac{1}{2} \mp \frac{iv}{\pi}\\
		0, & 0
	\end{matrix}\;\middle|\; 1 + e^{\pi\sigma}
	\right)
\end{align}
which further implies
\begin{align}
	\begin{split}
		&I_M^{\{\sech\}}(v) = 2 \sech(v) \\
		&\ - \frac{1}{\pi}\left(
		\frac{\left(e^{\pi\sigma}\right)^{\frac{1}{2}+\frac{iv}{\pi}}}{\Gamma\left(\frac{1}{2} + \frac{iv}{\pi}\right)}
		G_{2,2}^{2,2}
		\left(
		\begin{matrix}
			0, & \frac{1}{2} - \frac{iv}{\pi}\\
			0, & 0
		\end{matrix}\;\middle|\; 1 + e^{\pi\sigma}
		\right)
		-
		\frac{\left(e^{\pi\sigma}\right)^{\frac{1}{2}-\frac{iv}{\pi}}}{\Gamma\left(\frac{1}{2} - \frac{iv}{\pi}\right)}
		G_{2,2}^{2,2}
		\left(
		\begin{matrix}
			0, & \frac{1}{2} + \frac{iv}{\pi}\\
			0, & 0
		\end{matrix}\;\middle|\; 1 + e^{\pi\sigma}
		\right)
		\right).
	\end{split}
\end{align}

This form facilitates manipulation via Mathematica, which recognizes many alternate
representations of functions in Meijer-G form thus facilitating identifying a factorization
which isolates the oscillating factors.  We achieve this by identifying the form
\begin{align}
	\begin{split}
	I_M^{\{\sech\}}(v) &=
	\frac{\left(e^{\pi\sigma}\right)^{\frac{1}{2}+\frac{iv}{\pi}}}
	     {\left( e^{\pi\sigma} + 1 \right)\left( \frac{\pi}{2} - iv \right)}
	{}_2F_{1}\left(1, 1; \frac{3}{2} - \frac{iv}{\pi}; \frac{1}{1 + e^{\pi\sigma}} \right)
	\\
	&\quad\quad + \frac{\left(e^{\pi\sigma}\right)^{\frac{1}{2}-\frac{iv}{\pi}}}
	{\left( e^{\pi\sigma} + 1 \right)\left( \frac{\pi}{2} + iv \right)}
	{}_2F_{1}\left(1, 1; \frac{3}{2} + \frac{iv}{\pi}; \frac{1}{1 + e^{\pi\sigma}} \right).
	\end{split}
\end{align}
Noting that $\sigma$ increases with the number of terms $M$ and noting the limit
\begin{align}
	\lim_{\sigma \to \infty} {}_{2}F_{1}\left(1, 1; c; \frac{1}{1-e^{\pi \sigma}}\right) = 1
\end{align}
reveals the desired closed form expression
\begin{align}
	I_M^{\{\sech\}}(v) \simeq
	\frac{\left(e^{\pi\sigma}\right)^{\frac{1}{2}+\frac{iv}{\pi}}}
	{\left( e^{\pi\sigma} + 1 \right)\left( \frac{\pi}{2} - iv \right)}
	+ \frac{\left(e^{\pi\sigma}\right)^{\frac{1}{2}-\frac{iv}{\pi}}}
	{\left( e^{\pi\sigma} + 1 \right)\left( \frac{\pi}{2} + iv \right)}.
\end{align}
Simplifying trigonometrically yields
\begin{align}
	I_M^{\{\sech\}}(v) \simeq
	\frac{
		4 e^{\frac{\pi \sigma}{2}}
		\left(\pi \cos(\sigma v) - 2 v \sin(\sigma v)\right)
	}{(e^{\pi \sigma} + 1)(4v^2 + \pi^2)}.
\end{align}

\subsection{Hyperbolic Tangent}
We begin with the exact error formula as derived in Chapter \ref{chapter_improved_expansions}, $H_M^{\{\tanh\}}$, which follows directly from application of
(\ref{eq_general_remainder}):
\begin{align}
	\begin{split}
		&H_M^{\{\tanh\}}(v) = \int_\tau^{\infty} \!
		\csch\left(\frac{\pi}{2}\xi\right)
		\sin(\xi v) \mathrm{d}\xi
	\end{split}
	\\
	\begin{split}
		& \quad\quad + \frac{(-1)^{M+1}v^{2M+3}}{(2M+3)!} \int_0^{\tau} \!
		\xi^{2M+3} \csch\left(\frac{\pi}{2}\xi\right)
		{}_{1}F_{2}\left(1; M+2, M+\frac{5}{2}; -\frac{v^2 \xi^2}{4}\right) \mathrm{d}\xi.
	\end{split}
\end{align}

As noted in Chapter \ref{chapter_improved_expansions}, the second term dominates
for large $v$ due to the $v^{2M+3}$ factor, causing the expected divergence
for $|v| \geq V$.  Meanwhile, the first term (\ref{eq_tanh_remainder2})
dominates for small $v$ leading to Gibbs-like oscillation of the
error for $|v| < V$.  Noting this effect, the integral of
(\ref{eq_tanh_remainder2}) can be solved exactly to analyze the
error, $I_M^{\{\tanh\}}(v)$, inside the region of convergence for
an $M$ term expansion, i.e.,
\begin{gather}
	\begin{split}
		I_M^{\{\tanh\}}(v) &= \int_\tau^{\infty} \!
		\csch\left(\frac{\pi}{2}\xi\right)
		\sin(\xi v) \mathrm{d}\xi \\
		&= 2 \tanh(v) - \frac{i}{\pi}\left(
		\mathrm{B}_{e^{\pi \tau}}\left(\frac{1}{2} + \frac{iv}{\pi}, 0 \right)
		- \mathrm{B}_{e^{\pi \tau}}\left(\frac{1}{2} - \frac{iv}{\pi}, 0 \right)\right). \\
	\end{split}
\end{gather}
We will now reduce this result to a convenient approximate closed form which
captures most of the Gibbs-like oscillation.  We begin by noting a series of
equivalent representations of the integral.  The first of which employs the
identity
\begin{align}
	{}_{2}F_{1}\left(a; b; b+1; z\right) = bz^{-b}\mathrm{B}_z(b, 1-a)
\end{align}
from \cite{prudnikov_integrals_2003}
to make the substitution
\begin{align}
	\mathrm{B}_{e^{\pi\tau}}\left(\frac{1}{2} \pm \frac{iv}{\pi}, 0\right)
	= \frac{(e^{\pi\tau})^{\frac{1}{2} \pm \frac{iv}{\pi}}}{\frac{1}{2} \pm \frac{iv}{\pi}}
	{}_{2}F_{1}\left(1, \frac{1}{2} \pm \frac{iv}{\pi}; \frac{3}{2} \pm \frac{iv}{\pi}; e^{\pi \tau}\right).
\end{align}
This results in
\begin{align}
	I_M^{\{\tanh\}}(v) &= 2 \tanh(v) - \frac{i}{\pi}\left(u - u^*\right), \\
	u &= \frac{(e^{\pi\tau})^{\frac{1}{2} + \frac{iv}{\pi}}}{\frac{1}{2} + \frac{iv}{\pi}}
	    {}_{2}F_{1}\left(1, \frac{1}{2} + \frac{iv}{\pi}; \frac{3}{2} + \frac{iv}{\pi}; e^{\pi \tau}\right).
\end{align}
This can be equivalently expressed in terms of the Lerch Transcendent, $\Phi(z, a, b)$, by noting
from \cite{prudnikov_integrals_2003} that
\begin{align}
	{}_{2}F_{1}\left(1, b; b+1; z\right) = b \Phi(z, 1, b)
\end{align}
implies
\begin{align}
	{}_{2}F_{1}\left(1, \frac{1}{2} \pm \frac{iv}{\pi}; \frac{3}{2} \pm \frac{iv}{\pi}; e^{\pi \tau}\right)
	 = \left(\frac{1}{2} \pm \frac{iv}{\pi}\right)\Phi\left(e^{\pi\tau}, 1, \frac{1}{2} \pm \frac{iv}{\pi}\right)
\end{align}
which further implies
\begin{align}
\begin{split}
	&I_M^{\{\tanh\}}(v) = 2 \tanh(v) \\
	&\quad - \frac{i}{\pi}\left(
	    \left(e^{\pi\tau}\right)^{\frac{1}{2}+\frac{iv}{\pi}} \Phi \left(e^{\pi\tau}, 1, \frac{1}{2} + \frac{iv}{\pi} \right)
	    - \left(e^{\pi\tau}\right)^{\frac{1}{2}-\frac{iv}{\pi}} \Phi \left(e^{\pi\tau}, 1, \frac{1}{2} - \frac{iv}{\pi} \right)
	\right).
\end{split}
\end{align}
Noting that the identity
\begin{align}
	\Phi\left(1-z, 1, a\right)
	= \frac{1}{\Gamma(a)}
	G_{2,2}^{2,2}
	\left(
	\begin{matrix}
		0, & 1-a \\
		0, & 0
	\end{matrix}\;\middle|\;z
	\right)
\end{align}
from \cite{prudnikov_integrals_2003} implies
\begin{align}
	\Phi\left(e^{\pi\tau}, 1, \frac{1}{2} \pm \frac{iv}{\pi}\right)
	= \frac{1}{\Gamma\left(\frac{1}{2} \pm \frac{iv}{\pi}\right)}
	G_{2,2}^{2,2}
	\left(
	\begin{matrix}
		0, & \frac{1}{2} \mp \frac{iv}{\pi}\\
		0, & 0
	\end{matrix}\;\middle|\; 1 - e^{\pi\tau}
	\right)
\end{align}
which further implies
\begin{align}
	\begin{split}
	&I_M^{\{\tanh\}}(v) = 2 \tanh(v) \\
	&\ - \frac{i}{\pi}\left(
	\frac{\left(e^{\pi\tau}\right)^{\frac{1}{2}+\frac{iv}{\pi}}}{\Gamma\left(\frac{1}{2} + \frac{iv}{\pi}\right)}
	G_{2,2}^{2,2}
	\left(
	\begin{matrix}
		0, & \frac{1}{2} - \frac{iv}{\pi}\\
		0, & 0
	\end{matrix}\;\middle|\; 1 - e^{\pi\tau}
	\right)
	-
	\frac{\left(e^{\pi\tau}\right)^{\frac{1}{2}-\frac{iv}{\pi}}}{\Gamma\left(\frac{1}{2} - \frac{iv}{\pi}\right)}
	G_{2,2}^{2,2}
	\left(
	\begin{matrix}
		0, & \frac{1}{2} + \frac{iv}{\pi}\\
		0, & 0
	\end{matrix}\;\middle|\; 1 - e^{\pi\tau}
	\right)
	\right).
    \end{split}
\end{align}

This form facilitates manipulation via Mathematica, which recognizes many alternate
representations of functions in Meijer-G form thus facilitating identifying a factorization
which isolates the oscillating factors.  We achieve this by identifying the form
\begin{align}
\begin{split}
	I_M^{\{\tanh\}}(v) &= 2 \tanh(v) - \frac{i}{\pi}(u_1 - u_2), \\
    u_1 &= \frac{2\pi \left(e^{\pi\tau}\right)^{\frac{1}{2}+\frac{iv}{\pi}} {}_{2}F_{1}\left(1, 1; \frac{3}{2} - \frac{iv}{\pi}; \frac{1}{1-e^{\pi\tau}}\right)}{\left(e^{\pi\tau} - 1\right)(\pi - 2iv)} + \pi\left(-e^{\pi\tau}\right)^{-\frac{1}{2} - \frac{iv}{\pi}}\sech(v), \\
    u_2 &= \frac{2\pi \left(e^{\pi\tau}\right)^{\frac{1}{2}-\frac{iv}{\pi}} {}_{2}F_{1}\left(1, 1; \frac{3}{2} + \frac{iv}{\pi}; \frac{1}{1-e^{\pi\tau}}\right)}{\left(e^{\pi\tau} - 1\right)(\pi + 2iv)} + \pi\left(-e^{\pi\tau}\right)^{-\frac{1}{2} + \frac{iv}{\pi}}\sech(v).
\end{split}
\end{align}
Simplifying the terms yields
\begin{align}
	u_1 &=
	\frac{e^{\frac{\pi\tau}{2}}}{e^{\pi\tau} - 1}
	\frac{e^{i\pi\tau}}{\left(\frac{1}{2} - \frac{iv}{\pi}\right)}
	{}_{2}F_{1}\left(1, 1; \frac{3}{2} - \frac{iv}{\pi}; \frac{1}{1 - e^{\pi\tau}}\right)
	- i\pi (\tanh(v)+1), \\
	u_2 &=
	\frac{e^{\frac{\pi\tau}{2}}}{e^{\pi\tau} - 1}
	\frac{e^{-i\pi\tau}}{\left(\frac{1}{2} + \frac{iv}{\pi}\right)}
	{}_{2}F_{1}\left(1, 1; \frac{3}{2} + \frac{iv}{\pi}; \frac{1}{1 - e^{\pi\tau}}\right)
	- i\pi (1-\tanh(v)).
\end{align}
Further simplifying yields
\begin{align}
\begin{split}
	&I_M^{\{\tanh\}}(v) = \\
	&\quad \frac{i}{\pi}\frac{e^{\frac{\pi\tau}{2}}}{e^{\pi\tau} - 1}
	\frac{e^{-i\pi\tau}}{\left(\frac{1}{2} + \frac{iv}{\pi}\right)}
	{}_{2}F_{1}\left(1, 1; \frac{3}{2} + \frac{iv}{\pi}; \frac{1}{1 - e^{\pi\tau}}\right) \\
    &\quad\quad -
    \frac{i}{\pi}\frac{e^{\frac{\pi\tau}{2}}}{e^{\pi\tau} - 1}
	\frac{e^{i\pi\tau}}{\left(\frac{1}{2} - \frac{iv}{\pi}\right)}
	{}_{2}F_{1}\left(1, 1; \frac{3}{2} - \frac{iv}{\pi}; \frac{1}{1 - e^{\pi\tau}}\right)
\end{split}
\end{align}
which can in turn be simplified to an exact expression for $I_M^{\{\tanh\}}(v)$, i.e.,
\begin{align}
	I_M^{\{\tanh\}}(v) &= 
	\frac{e^{\frac{\pi\tau}{2}}}{e^{\pi\tau} - 1}\left(u - u*\right), \\
	u &= \frac{e^{-i\tau v}}{v - i\frac{\pi}{2}}
	     {}_{2}F_{1}\left(1, 1; \frac{3}{2} + \frac{iv}{\pi}; \frac{1}{1 - e^{\pi\tau}}\right).
\end{align}
Noting that $\tau$ increases with the number of terms $M$ and noting the limit
\begin{align}
	\lim_{\tau \to \infty} {}_{2}F_{1}\left(1, 1; c; \frac{1}{1-e^{\pi \tau}}\right) = 1
\end{align}
reveals the desired closed form expression
\begin{align}
	I_M^{\{\tanh\}}(v) &\simeq 
	\frac{e^{\frac{\pi\tau}{2}}}{e^{\pi\tau} - 1}\left(
	\frac{e^{-i\tau v}}{v - \frac{i\pi}{2}} + \frac{e^{i\tau v}}{v + \frac{i \pi}{2}}
	\right).
\end{align}
Simplifying trigonometrically yields
\begin{align}
	I_M^{\{\tanh\}}(v) \simeq
	\frac{
		4 e^{\frac{\pi \tau}{2}}
		\left(\pi \sin(\tau v) + 2 v \cos(\tau v)\right)
	}{(e^{\pi \tau} - 1)(4v^2 + \pi^2)}.
\end{align}
